%! AP Statistics — Unit 4, Lesson 4.8 — Student Packet Cover Sheet
\documentclass[10pt]{article}
\usepackage{apstats-article}
\usepackage{apstats-boxes}
\usepackage{ltablex}
\keepXColumns

\begin{document}

% Full-bleed navy banner
\begin{tikzpicture}[remember picture, overlay]
  \fill[navy] (current page.north west)
    rectangle ([yshift=-0.9in]current page.north east);
\end{tikzpicture}
\vspace{-0.8in}

\begin{center}
  {\color{white}\textbf{\LARGE AP Statistics}}\\[4pt]
  {\color{skymid}\large\textbf{Unit 4: Probability, Random Variables, and Probability Distributions}}\\[6pt]
  {\color{white}\textbf{\large Lesson 4.8 \quad Mean and Standard Deviation of Random Variables}}
\end{center}

\vspace{0.22in}
\namedateperiod
\vspace{0.18in}

\begin{learningtargetbox}
\small
\begin{enumerate}[leftmargin=1.5em, itemsep=5pt]
  \item \ldots{} compute the mean (expected value) of a discrete random variable using
        the formula $\mu_X = \sum x \cdot P(X=x)$.
  \item \ldots{} compute the standard deviation of a discrete random variable using
        $\sigma_X = \sqrt{\sum (x-\mu_X)^2 \cdot P(X=x)}$.
  \item \ldots{} interpret $\mu_X$ and $\sigma_X$ in the context of the random variable.
\end{enumerate}
\end{learningtargetbox}

\vspace{0.14in}

\begin{tocbox}
\small
\renewcommand{\arraystretch}{1.5}
\begin{tabularx}{\linewidth}{@{} c l X r @{}}
\rowcolor{sky}
\textbf{\#} & \textbf{Component} & \textbf{Description} & \textbf{Score} \\
\midrule
1 & Warm-Up        & Spiral review: distributions from Lesson 4.7          & \blank{1.2cm} \\
2 & Guided Notes   & Computation table; $\mu_X$ and $\sigma_X$ formulas    & \blank{1.2cm} \\
3 & Group Activity & Computing and interpreting $\mu_X$ and $\sigma_X$     & \blank{1.2cm} \\
4 & Exit Ticket    & Goals distribution; full computation                  & \blank{1.2cm} \\
5 & Homework       & Expected value and SD; contextual interpretations      & \blank{1.2cm} \\
\midrule
  & \multicolumn{2}{r}{\textbf{Total}} & \blank{1.2cm} \\
\end{tabularx}
\end{tocbox}

\vspace{0.14in}

\begin{remindbox}
\small
The \textbf{most important idea in Lesson 4.8}: the expected value is a long-run average,
not a prediction for any single outcome.

\vspace{6pt}
\begin{tabularx}{\linewidth}{@{} >{\centering\arraybackslash\columncolor{goldbg}}p{0.06\linewidth}
                                    X @{}}
\textbf{\large!} &
\textbf{$\mu_X$ need not be a possible value of $X$.}\enspace
For example, $\mu_X = 2.3$ children per family is meaningful even though no family
has 2.3 children. \\[6pt]
\textbf{\large!} &
\textbf{Variances add; standard deviations do not.}\enspace
Coming up in Lesson 4.9: when combining random variables, always add variances first,
then take the square root. \\[6pt]
\textbf{\large!} &
\textbf{Interpretation earns AP points.}\enspace
Always say what $\mu_X$ and $\sigma_X$ mean for the specific situation --- do not
just state the numbers. \\
\end{tabularx}
\end{remindbox}

\end{document}
